% ============================================================================
% ARBEITSBLATT DS 3 - Gruppe A: Modalverben (S. 32-33)
% ============================================================================

\documentclass[11pt,a4paper]{article}

\usepackage[utf8]{inputenc}
\usepackage[T1]{fontenc}
\usepackage[ngerman]{babel}
\usepackage[left=2cm,right=2cm,top=1.8cm,bottom=1.5cm]{geometry}
\usepackage{helvet}
\usepackage[table]{xcolor}
\usepackage{tabularx}
\usepackage{array}
\usepackage{enumitem}
\usepackage{tcolorbox}
\usepackage{fancyhdr}
\usepackage{lastpage}

\renewcommand{\familydefault}{\sfdefault}

\definecolor{spanisch}{HTML}{c41e3a}
\definecolor{grau}{HTML}{666666}
\definecolor{hellgrau}{HTML}{f5f5f5}

\pagestyle{fancy}
\fancyhf{}
\fancyhead[L]{\small\textcolor{grau}{Spanisch Spätbeginner -- Gruppe A}}
\fancyhead[R]{\small\textcolor{grau}{AB-03-A}}
\fancyfoot[C]{\small\thepage/\pageref{LastPage}}
\renewcommand{\headrulewidth}{0.5pt}

\newcommand{\luecke}[1]{\underline{\hspace{#1}}}
\newcommand{\punktebox}[1]{\hfill\fbox{\small \_\_/#1}}
\newcommand{\es}[1]{\textit{#1}}

\newtcolorbox{merkkasten}[1][]{
    colback=white,
    colframe=spanisch,
    fonttitle=\bfseries\color{spanisch},
    title=#1,
    boxrule=1.5pt,
    arc=0pt,
    left=5pt, right=5pt, top=5pt, bottom=5pt
}

\newtcolorbox{buchverweis}{
    colback=hellgrau,
    colframe=grau,
    boxrule=0.5pt,
    arc=0pt,
    left=5pt, right=5pt, top=3pt, bottom=3pt
}

\newcounter{aufgabe}
\newcommand{\aufgabe}[2]{%
    \stepcounter{aufgabe}%
    \vspace{0.8em}%
    \noindent\textbf{\theaufgabe.} #1 \punktebox{#2}%
    \vspace{0.3em}%
}

\begin{document}

% ============================================================================
% KOPFBEREICH
% ============================================================================
\noindent
\begin{tabularx}{\textwidth}{@{}Xr@{}}
    {\Large\textbf{\textcolor{spanisch}{Poder, querer, tener que}}} & Name: \luecke{5cm} \\[0.3em]
    {\small DS 3 -- Modalverben (S. 32--33)} & Datum: \luecke{3cm} \\
\end{tabularx}

\vspace{0.3em}
\hrule
\vspace{0.8em}

% ============================================================================
% MERKKASTEN 1: Modalverben
% ============================================================================
\begin{merkkasten}[Kasten 1: Modalverben (können / wollen / müssen)]
\vspace{0.2em}

\begin{tabularx}{\textwidth}{@{}|c|l|l|l|@{}}
\hline
\rowcolor{hellgrau}
& \textbf{poder} (können) & \textbf{querer} (wollen) & \textbf{tener que} (müssen) \\
\hline
yo & \luecke{2cm} & \luecke{2cm} & tengo que \\
\hline
tú & puedes & \luecke{2cm} & tienes que \\
\hline
él/ella & \luecke{2cm} & quiere & \luecke{2cm} \\
\hline
nosotros & podemos & \luecke{2cm} & tenemos que \\
\hline
vosotros & \luecke{2cm} & queréis & \luecke{2cm} \\
\hline
ellos & pueden & \luecke{2cm} & tienen que \\
\hline
\end{tabularx}

\vspace{0.3em}
\textbf{Struktur:} Modalverb + Infinitiv \quad Beispiel: \es{Puedo ir} al cine.
\vspace{0.2em}
\end{merkkasten}

\vspace{1em}

% ============================================================================
% AUFGABEN
% ============================================================================

\aufgabe{Ergänze die Konjugationstabelle oben.}{9}

\vspace{0.3em}

\aufgabe{Ergänze die richtige Form von \es{poder}, \es{querer} oder \es{tener que}.}{8}

\begin{enumerate}[label=\alph*), leftmargin=2em]
    \item ¿\luecke{2.5cm} (tú/poder) ir al cine esta noche?
    \item Nosotros \luecke{2.5cm} (querer) ir a la piscina.
    \item Ellos \luecke{2.5cm} (tener que) estudiar para el examen.
    \item Yo no \luecke{2.5cm} (poder) salir hoy.
    \item ¿Qué \luecke{2.5cm} (vosotros/querer) hacer el sábado?
    \item María \luecke{2.5cm} (tener que) trabajar el domingo.
    \item ¿\luecke{2.5cm} (nosotros/poder) ir al parque después?
    \item Él \luecke{2.5cm} (querer) aprender español.
\end{enumerate}

\vspace{0.5em}

\aufgabe{Was bedeutet was? Verbinde. (Buch S. 32, Aufg. 4)}{4}

\begin{tabularx}{\textwidth}{@{}|X|X|@{}}
\hline
\rowcolor{hellgrau}
\textbf{Spanisch} & \textbf{Deutsch} \\
\hline
1. No puedo ir porque tengo que estudiar. & a) Willst du ins Kino gehen? \\
\hline
2. ¿Quieres ir al cine? & b) Wir müssen um 8 zu Hause sein. \\
\hline
3. Tenemos que estar en casa a las 8. & c) Ich kann nicht gehen, weil ich lernen muss. \\
\hline
4. ¿Podéis venir a mi fiesta? & d) Könnt ihr zu meiner Party kommen? \\
\hline
\end{tabularx}

\vspace{0.3em}
Lösungen: 1 = \luecke{0.5cm}, 2 = \luecke{0.5cm}, 3 = \luecke{0.5cm}, 4 = \luecke{0.5cm}

\vspace{0.8em}

\aufgabe{Schreibe Sätze mit Modalverben. (Buch S. 33, Aufg. 5)}{6}

\begin{enumerate}[label=\alph*), leftmargin=2em]
    \item (yo / querer / ir / al cine)

    \luecke{12cm}

    \item (nosotros / poder / jugar / al fútbol)

    \luecke{12cm}

    \item (ella / tener que / hacer / los deberes)

    \luecke{12cm}
\end{enumerate}

\vspace{0.8em}

\aufgabe{Dialog: Wochenendpläne mit Modalverben.}{8}

\begin{buchverweis}
\textbf{Schreibe einen Dialog (6 Sätze).} Verwendet: \es{querer}, \es{poder}, \es{tener que}.
\end{buchverweis}

\vspace{0.3em}
\begin{tabularx}{\textwidth}{@{}|X|@{}}
\hline
\\[0.4em]
\textbf{A:} \luecke{13cm} \\[0.7em]
\textbf{B:} \luecke{13cm} \\[0.7em]
\textbf{A:} \luecke{13cm} \\[0.7em]
\textbf{B:} \luecke{13cm} \\[0.7em]
\textbf{A:} \luecke{13cm} \\[0.7em]
\textbf{B:} \luecke{13cm} \\[0.4em]
\hline
\end{tabularx}

\vspace{1em}
\hrule
\vspace{0.3em}
\begin{flushright}
\textbf{Gesamt: \luecke{1cm} / 35 Punkte}
\end{flushright}

\end{document}
